\documentclass[12pt,a4paper,oneside]{report}
\usepackage[T1]{fontenc}
\usepackage[utf8]{inputenc}
\usepackage[spanish,es-tabla]{babel}
\usepackage{lmodern}
\usepackage{geometry}
\geometry{margin=2.5cm}
\usepackage{fancyhdr}
\setlength{\headheight}{15pt}
\pagestyle{fancy}
\fancyhf{}
\fancyhead[L]{\small\scshape\nouppercase{\leftmark}}
\fancyfoot[C]{\thepage}
\renewcommand{\headrulewidth}{0.4pt}
\usepackage{setspace}
\onehalfspacing
\usepackage{graphicx}
\graphicspath{{figuras/}}
\usepackage{amsmath,amssymb,booktabs}
\usepackage[hidelinks]{hyperref}
\newcommand{\titulotfm}{Optimización de modelos generativos avanzados para la síntesis de imágenes agrícolas orientadas a la detección automática de \textit{Amaranthus palmeri}}
\newcommand{\autortfm}{Christian Andrés Rueda Ayala}
\newcommand{\tutortfm}{Ph.D. Raquel Martínez}
\newcommand{\mastertfm}{Máster Universitario en Inteligencia Artificial}
\newcommand{\cursotfm}{2026--2027}
\begin{document}
\hypersetup{pageanchor=false}
\begin{titlepage}
\centering
\vspace*{0.5cm}
% Logo de la Universidad de Murcia guardado en doc/recursos.
\includegraphics[width=7.5cm]{recursos/02_umu.png}\par
\vspace{1.2cm}
{\large\scshape Universidad de Murcia\par}
\vspace{0.15cm}
{\scshape Trabajo Fin de M\'aster\par}
{\small\mastertfm\par}
\vspace{0.5cm}
\noindent\rule{\textwidth}{0.4pt}\par
\vspace{0.2cm}
{\Large\bfseries \titulotfm\par}
\vspace{0.2cm}
\noindent\rule{\textwidth}{0.4pt}\par
\vspace{0.9cm}
\begin{minipage}[t]{0.46\textwidth}
\raggedright\textit{Autor/a:}\\\textsc{\autortfm}
\end{minipage}\hfill
\begin{minipage}[t]{0.46\textwidth}
\raggedleft\textit{Tutor/a:}\\\textsc{\tutortfm}
\end{minipage}
\vfill
{\small \cursotfm\par}
\end{titlepage}

\hypersetup{pageanchor=true}
\pagenumbering{roman}
% Opcional: descomenta para incluir agradecimientos.
% \chapter*{Agradecimientos}
% Escribe aquí tus agradecimientos.

\tableofcontents
\chapter*{Resumen}
\addcontentsline{toc}{chapter}{Resumen}
Escribe aquí el resumen del trabajo.

\noindent\textbf{Palabras clave:} inteligencia artificial; añadir palabras clave.

\chapter*{Abstract}
\addcontentsline{toc}{chapter}{Abstract}
Write the abstract here.

\noindent\textbf{Keywords:} artificial intelligence; add keywords.

\clearpage
\pagenumbering{arabic}
\chapter{Introducción}
\section{Contexto y motivación}
\section{Objetivos}
\section{Estructura de la memoria}

\chapter{Metodología}

Este capítulo describe el diseño del estudio y la preparación del corpus disponible hasta el 12 de noviembre de 2026. Se distinguen los procedimientos realizados de las fases experimentales previstas. El objetivo es disponer de imágenes agrícolas trazables para estudiar la síntesis y validación de imágenes de \textit{Amaranthus palmeri}, combinando criterios computacionales con la revisión de su identidad y morfología vegetal.

\section{Diseño general del estudio}

El estudio parte de imágenes reales de tres especies del género \textit{Amaranthus}: \textit{A. palmeri}, \textit{A. hybridus} y \textit{A. retroflexus}. Las tres se consideran en el desarrollo previsto de un clasificador, mientras que la generación se centrará en \textit{A. palmeri}. Las imágenes de maíz (\textit{Zea mays}) se conservan como material auxiliar, separado de las clases principales y de los ejemplos positivos destinados al modelo generativo. Para mantener juntas las observaciones relacionadas, la unidad de agrupación se define por año, campaña, repetición y posición de la maceta.

El flujo previsto comprende la consolidación y revisión de los datos, la asignación de reservas por grupos, la preparación de imágenes con una referencia física común, el desarrollo de un clasificador y la adaptación de un modelo generativo. Posteriormente se prevé auditar las imágenes generadas mediante el clasificador, representaciones de BioCLIP, evaluación experta y controles geométricos. En la fecha de corte se han consolidado los manifiestos y comprobado su integridad; las particiones, el preprocesado y los experimentos de modelado permanecen pendientes.

\section{Material vegetal y adquisición de imágenes}

El corpus actual reúne 1.451 fotografías JPEG correspondientes a 15 sesiones de las campañas de 2025 y 2026. Los manifiestos consolidados comprenden 493 fotografías de 2025 y 958 de 2026. Se trata de fotografías de macetas adquiridas en distintas fechas y, en algunos casos, con cámaras diferentes. Según el investigador responsable de la adquisición, dentro de cada campaña y repetición las fotografías de una posición corresponden a las mismas macetas y plantas a lo largo del tiempo. Esta continuidad permite relacionar observaciones longitudinales, pero no demuestra por sí sola la independencia entre repeticiones ni descarta conexiones por trasplante.

El diseño de siembra se documentó en cuadrículas con recorrido en serpentina. En 2025 se registró para cada posición el número de plantas sembradas; este dato se vinculó al diseño siguiendo el orden del recorrido. En 2026 se registró una planta nominal por maceta según la declaración del investigador. Se conserva la distinción entre el número sembrado o nominal y el número de plantas efectivamente visibles o vivas en una fotografía, que no se ha contabilizado de forma exhaustiva.

Para la futura evaluación geométrica se dispone del diámetro exterior nominal de la boca de la maceta: 10 cm para 2025 y 20 cm para 2026. Estos valores son referencias documentales; todavía no se han medido los diámetros en píxeles ni se ha validado una escala para cada imagen. La presencia de tarjetas se conserva como información esperada según las notas de campo, sin tratarla como una observación exhaustiva ni como referencia de escala adoptada.

\section{Inventario y comprobación de integridad}

Se realizó un inventario recursivo de las imágenes y sus archivos auxiliares. La auditoría inicial incluyó la decodificación de los JPEG, la extracción de metadatos EXIF, incluidos los contenidos en directorios anidados, y el cálculo de SHA-256 de archivo y de hashes de píxeles RGB tras aplicar la orientación EXIF. En la instantánea inicial se identificaron 1.474 fotografías, sin fallos de lectura ni duplicados exactos según las comprobaciones efectuadas. Posteriormente, el investigador retiró 23 fotografías de una secuencia repetida; el corpus actual consta de 1.451 archivos. Los recuentos de ambas fechas se conservan separados para mantener la trazabilidad del cambio.

En una fotografía Nikon de la sesión del 20 de junio de 2025 se detectó un cambio tanto en el hash del archivo como en el de los píxeles orientados respecto de la auditoría inicial. Se conservó y verificó la versión disponible en el momento de la consolidación. La causa del cambio y su relevancia visual no se han establecido; su revisión permanece pendiente.

\section{Asignación de posiciones y etiquetas}

Las fotografías se relacionaron con el diseño mediante la posición en el recorrido en serpentina y el orden temporal de adquisición. La fecha de sesión se tomó del nombre de la carpeta, mientras que las fechas y horas EXIF originales se conservaron por separado. Esta decisión atiende a una discrepancia observada entre la fecha de carpeta y la fecha registrada por la Nikon D70, sin modificar los metadatos de origen ni atribuir de forma concluyente la causa del desfase. Cuando la información temporal no permitiera resolver una secuencia, un desempate por nombre no se consideraría una confirmación del orden real.

Las aclaraciones del investigador permitieron resolver casos específicos. En la sesión del 1 de junio de 2026 se aplicó el orden cronológico confirmado, comenzando por el bloque de archivos 0856--0953 y continuando por el bloque 0008--0043. Para el 16 de junio de 2025 se trataron las 60 fotografías disponibles como una selección de una toma por posición de \textit{Amaranthus}, y no como parejas de tomas de 30 posiciones. Para el 20 de junio de 2025 se asignaron por separado las adquisiciones de iPhone y Nikon: las 23 fotografías Nikon se vincularon a las primeras 23 posiciones cubiertas también por iPhone, sin contarlas como macetas adicionales. La sesión incompleta se mantuvo como tal, sin inventar observaciones de las posiciones no fotografiadas.

La especie prevista en el diseño se conservó separada de la especie de trabajo. Tras revisar las posiciones F26C1 y F4C4 de la segunda repetición de 2026, el investigador identificó sus plantas como \textit{A. palmeri}, aunque el diseño preveía \textit{A. hybridus}. La corrección se aplicó a las ocho fotografías disponibles de esas dos macetas y se registraron su fuente y alcance, conservando la etiqueta original. Esta identificación procede de la revisión del investigador; no se presenta como una determinación molecular ni como una validación por un panel independiente. Las demás etiquetas de diseño no se consideran confirmadas individualmente por el mero hecho de figurar en el manifiesto.

\section{Copias de revisión y manifiestos consolidados}

Se preparó una copia de cada fotografía actual para facilitar la revisión visual de la siembra. Los nombres de estas copias incorporan el orden asignado, la especie esperada según el diseño, la posición y el nombre original. Se verificó mediante SHA-256 que las 1.451 copias coincidían con los archivos de origen disponibles en ese momento. Los nombres conservan la expectativa inicial del diseño y no constituyen la fuente de la etiqueta de trabajo corregida. Se registraron las ocho observaciones correspondientes a las dos posiciones corregidas; la preparación de las copias no implica la revisión individual del resto del corpus.

Los manifiestos por año y el manifiesto conjunto se construyeron con \texttt{scripts/build\_dataset\_manifests.py}, incorporando las correcciones explícitas de \texttt{exp\_design/species\_corrections.csv}. Para cada fotografía se conservaron, entre otros campos, la procedencia, la sesión, la cámara, la posición, las especies de diseño y de trabajo, el grupo de maceta, la información de siembra, las incidencias, los metadatos de imagen y los hashes. Los datos aún no observados o calculados se dejaron sin completar. Los consolidados anteriores se preservaron como versiones históricas.

La validación independiente de los manifiestos incluyó 54 comprobaciones satisfactorias, entre ellas la cobertura de los originales, la correspondencia por SHA-256, la coherencia de las correcciones y grupos y la equivalencia entre las salidas CSV y JSON. Esta validación verifica la coherencia técnica y documental de los archivos; no sustituye la revisión botánica de cada planta.

\section{Criterios preliminares de selección}

Se registraron nueve imágenes excluidas del entrenamiento principal: una con una banda o discontinuidad visual, tres anotadas sin planta y cinco anotadas con plantas secas o muertas. Estas imágenes se conservan en los manifiestos y en las copias de revisión, con el motivo correspondiente. Una planta pequeña no se excluye únicamente por su tamaño. La elegibilidad de los demás archivos sigue sujeta a la revisión de contenido y a los criterios de preparación que se definan.

Con las etiquetas de trabajo y exclusiones vigentes, se identificaron 1.319 imágenes candidatas para clasificación, correspondientes a 360 grupos: 440 de \textit{A. palmeri}, 434 de \textit{A. hybridus} y 445 de \textit{A. retroflexus}. También se identificaron 440 candidatas de \textit{A. palmeri} para una adaptación mediante LoRA. Estos dos conjuntos se solapan y no son reservas independientes ni tamaños finales de entrenamiento. El material auxiliar de maíz comprende 124 fotografías; una de ellas también figura entre las nueve exclusiones, por lo que ambas cantidades no deben sumarse como categorías disjuntas.

\section{Fases experimentales pendientes}

Antes de entrenar los modelos se prevé asignar las reservas de entrenamiento y validación del clasificador, entrenamiento del generador y prueba final por grupos de maceta, manteniendo el maíz como conjunto auxiliar. La asignación deberá comprobar la ausencia de grupos compartidos y documentar las limitaciones de independencia entre repeticiones. En la fecha de corte, todas las particiones permanecen sin asignar.

El preprocesado previsto incluirá la orientación de las imágenes, la medición del borde exterior de la maceta y la elección de un campo físico común a partir de datos de desarrollo. El diámetro nominal no equivale a una medición de escala por imagen. Tampoco se han generado maestros, fijado el campo común ni validado las proporciones. El recuento de plantas observadas, la visibilidad, el estado y el estadio se revisarán cuando sean necesarios para los lotes seleccionados.

Como arquitectura de clasificación se prevé explorar ConvNeXt-Tiny y utilizar BioCLIP como instrumento complementario de auditoría. Para la generación se han previsto tres adaptaciones LoRA con configuración común y semillas diferentes. La posibilidad de generar la planta sobre fondos de maceta y suelo se mantiene como objetivo de investigación: todavía deben definirse y probarse el procedimiento y las regiones necesarias. El uso de escenas completas no garantiza que el modelo aprenda exclusivamente la planta. No se han preparado las leyendas de entrenamiento ni se han ejecutado entrenamientos, generaciones o evaluaciones de modelos y expertos; sus configuraciones finales y resultados se documentarán cuando estén disponibles.

\chapter{Resultados}

\chapter{Discusión}

\chapter{Conclusiones}

\bibliographystyle{plain}
\bibliography{bibliografia/01_referencias}
\appendix
\chapter{Material complementario}

\end{document}
